% CuTe 白皮书参考文献：按官方英文 PDF 的 29 条及原顺序补齐。
% 使用显式编号段落，保证 XeLaTeX 与 Pandoc/MyST 均保留编号。
\section*{参考文献}
\addcontentsline{toc}{section}{参考文献}
\markboth{参考文献}{参考文献}

\noindent\textbf{[1]} NVIDIA. \emph{PyCuTe: A Pure-Python Reference Implementation of CuTe}, 2026.
\url{https://github.com/NVlabs/CuTe}. Accessed: 2026-07-28.

\noindent\textbf{[2]} NVIDIA. \href{https://images.nvidia.com/content/volta-architecture/pdf/volta-architecture-whitepaper.pdf}{\emph{Volta Architecture Whitepaper}}, 2017.

\noindent\textbf{[3]} NVIDIA. \href{https://images.nvidia.com/aem-dam/en-zz/Solutions/design-visualization/technologies/turing-architecture/NVIDIA-Turing-Architecture-Whitepaper.pdf}{\emph{Turing Architecture Whitepaper}}, 2018.

\noindent\textbf{[4]} NVIDIA. \href{https://images.nvidia.com/aem-dam/en-zz/Solutions/data-center/nvidia-ampere-architecture-whitepaper.pdf}{\emph{Ampere Architecture Whitepaper}}, 2021.

\noindent\textbf{[5]} NVIDIA. \href{https://resources.nvidia.com/en-us-hopper-architecture/nvidia-h100-tensor-c}{\emph{Hopper Architecture Whitepaper}}, 2023.

\noindent\textbf{[6]} NVIDIA. \href{https://images.nvidia.com/aem-dam/Solutions/geforce/blackwell/nvidia-rtx-blackwell-gpu-architecture.pdf}{\emph{Blackwell Architecture Whitepaper}}, 2025.

\noindent\textbf{[7]} L. Susan Blackford, Antoine Petitet, Roldan Pozo, Karin Remington,
R. Clint Whaley, James Demmel, Jack Dongarra, Iain Duff, Sven Hammarling,
Greg Henry, et al. An updated set of basic linear algebra subprograms (BLAS).
\emph{ACM Transactions on Mathematical Software}, 28(2):135--151, 2002.

\noindent\textbf{[8]} Field G. Van Zee and Robert A. van de Geijn. BLIS: A framework for rapidly
instantiating BLAS functionality. \emph{ACM Transactions on Mathematical Software},
41(3), 2015. \url{https://doi.org/10.1145/2764454}.

\noindent\textbf{[9]} Yang Shi, U. N. Niranjan, Animashree Anandkumar, and Cris Cecka.
Tensor Contractions with Extended BLAS Kernels on CPU and GPU.
\emph{2016 IEEE 23rd International Conference on High Performance Computing (HiPC)},
193--202, 2016. \url{https://doi.org/10.1109/HiPC.2016.031}.

\noindent\textbf{[10]} Ehud Aharoni, Allon Adir, Moran Baruch, Nir Drucker, Gilad Ezov,
Ariel Farkash, Lev Greenberg, Ramy Masalha, Guy Moshkowich, Dov Murik,
Hayim Shaul, and Omri Soceanu. HeLayers: A tile tensors framework for large
neural networks on encrypted data. \emph{Proceedings on Privacy Enhancing
Technologies}, 2023(1):325--342, 2023.
\url{https://doi.org/10.56553/popets-2023-0020}.

\noindent\textbf{[11]} Benjamin F. Spector, Simran Arora, Aaryan Singhal, Daniel Y. Fu, and
Christopher R\'e. ThunderKittens: Simple, fast, and adorable AI kernels, 2024.
\url{https://arxiv.org/abs/2410.20399}.

\noindent\textbf{[12]} Keren Zhou, Mario Lezcano, Adam Goucher, Akhmed Rakhmati, Jeff Niu,
Justin Lebar, Pawel Szczerbuk, Peter Bell, Phil Tillet, Thomas Raoux,
and Zahi Moudallal. Linear Layouts: Robust code generation of efficient
tensor computation using $\mathbb{F}_2$, 2025.
\url{https://arxiv.org/abs/2505.23819}.

\noindent\textbf{[13]} Philippe Tillet, Hsiang-Tsung Kung, and David D. Cox. Triton: An intermediate
language and compiler for tiled neural network computations.
\emph{Proceedings of the 3rd ACM SIGPLAN International Workshop on Machine
Learning and Programming Languages}, 2019.

\noindent\textbf{[14]} A. Edelman, S. Heller, and S. Lennart Johnsson. Index transformation algorithms
in a linear algebra framework. \emph{IEEE Transactions on Parallel and
Distributed Systems}, 5(12):1302--1309, 1994.
\url{https://doi.org/10.1109/71.334903}.

\noindent\textbf{[15]} T. H. Cormen. Fast permuting on disk arrays. \emph{Journal of Parallel and
Distributed Computing}, 17(1):41--57, 1993.
\url{https://doi.org/10.1006/jpdc.1993.1004}.

\noindent\textbf{[16]} Mathis Bouverot-Dupuis and Mary Sheeran. Efficient GPU implementation of affine
index permutations on arrays. \emph{Proceedings of the 11th ACM SIGPLAN
International Workshop on Functional High-Performance and Numerical Computing},
FHPNC 2023, 15--28, 2023. \url{https://doi.org/10.1145/3609024.3609411}.

\noindent\textbf{[17]} NVIDIA. \emph{CUTLASS: CUDA Templates for Linear Algebra Subroutines and Solvers},
2023. \url{https://github.com/NVIDIA/cutlass/tree/v3.0.0}.

\noindent\textbf{[18]} Somashekaracharya G. Bhaskaracharya, Aravind Acharya, Bastian Hagedorn,
and Vinod Grover. Modeling Layout Abstractions Using Integer Set Relations, 2025.
\url{https://arxiv.org/abs/2511.10374}.

\noindent\textbf{[19]} Amir Mohammad Tavakkoli, Cosmin Oancea, and Mary Hall. LEGO: Layout Expression
for Generating One-to-one Mapping, 2025. \url{https://arxiv.org/abs/2505.08091}.

\noindent\textbf{[20]} Jack Carlisle, Jay Shah, Reuben Stern, and Paul VanKoughnett. Categorical
Foundations for CuTe Layouts, 2026. \url{https://arxiv.org/abs/2601.05972}.

\noindent\textbf{[21]} NVIDIA. \href{https://docs.nvidia.com/cutlass/latest/media/docs/pythonDSL/cute_dsl.html}{\emph{CuTe DSL: Python APIs for CUTLASS}}, 2025. Accessed: 2026-03-01.

\noindent\textbf{[22]} Bastian Hagedorn, Bin Fan, Hanfeng Chen, Cris Cecka, Michael Garland,
and Vinod Grover. Graphene: An IR for optimized tensor computations on GPUs.
\emph{Proceedings of the 28th ACM International Conference on Architectural
Support for Programming Languages and Operating Systems, Volume 3}, ASPLOS 2023,
302--313, 2023. \url{https://doi.org/10.1145/3582016.3582018}.

\noindent\textbf{[23]} Muhammad Osama, Duane Merrill, Cris Cecka, Michael Garland, and John D. Owens.
Stream-K: Work-Centric Parallel Decomposition for Dense Matrix-Matrix
Multiplication on the GPU. \emph{Proceedings of the 28th ACM SIGPLAN Annual
Symposium on Principles and Practice of Parallel Programming}, PPoPP '23,
429--431, 2023. \url{https://doi.org/10.1145/3572848.3577479}.

\noindent\textbf{[24]} Andrew Kerr, Duane Merrill, Julien Demouth, and John Tran.
\href{https://developer.nvidia.com/blog/cutlass-linear-algebra-cuda}{\emph{CUTLASS: Fast Linear Algebra in CUDA C++}}, 2017.

\noindent\textbf{[25]} Tri Dao, Dan Fu, Stefano Ermon, Atri Rudra, and Christopher R\'e.
FlashAttention: Fast and Memory-Efficient Exact Attention with IO-Awareness.
\emph{Advances in Neural Information Processing Systems}, 35:16344--16359, 2022.

\noindent\textbf{[26]} Tri Dao. FlashAttention-2: Faster Attention with Better Parallelism and Work
Partitioning, 2023. \url{https://arxiv.org/abs/2307.08691}.

\noindent\textbf{[27]} Jay Shah, Ganesh Bikshandi, Ying Zhang, Vijay Thakkar, Pradeep Ramani,
and Tri Dao. FlashAttention-3: Fast and accurate attention with asynchrony and
low-precision. \emph{Advances in Neural Information Processing Systems},
37:68658--68685, 2024. \url{https://doi.org/10.52202/079017-2193}.

\noindent\textbf{[28]} Sharan Chetlur, Cliff Woolley, Philippe Vandermersch, Jonathan Cohen,
John Tran, Bryan Catanzaro, and Evan Shelhamer. cuDNN: Efficient Primitives
for Deep Learning, 2014. \url{https://arxiv.org/abs/1410.0759}.

\noindent\textbf{[29]} NVIDIA Corporation. \emph{PTX ISA 9.0 Documentation}. Technical report,
NVIDIA Corporation, 2025.
\url{https://docs.nvidia.com/cuda/parallel-thread-execution/index.html}.
